Considera la funció
\[
f(x)=x^4-4x^3 .
\]

\begin{apartats}

\apartat[1,25]{1}
Estudia la monotonia de $f$ i troba'n els extrems relatius.

\begin{solucio}
$f'(x)=4x^3-12x^2=4x^2(x-3)$, que s'anul·la en $x=0$ i en $x=3$.\\
Com que $4x^2\ge0$, el signe de $f'$ és el de $x-3$: negatiu a $(-\infty,3)$ i positiu a
$(3,+\infty)$.\\
Per tant, $f$ \textbf{decreix} a $(-\infty,3)$ i \textbf{creix} a $(3,+\infty)$.\\
En $x=3$ passa de decréixer a créixer: \textbf{mínim relatiu} $(3,-27)$.\\
En $x=0$ la derivada s'anul·la però \textbf{no canvia de signe}: no hi ha extrem. Que
$f'(a)=0$ no vol dir que en $a$ hi hagi un màxim o un mínim.
\end{solucio}

\begin{nomesllarg}
\apartat{0,75}
Aplica el criteri de la derivada segona als dos punts on s'anul·la $f'$. Què hi passa?

\begin{solucio}
$f''(x)=12x^2-24x=12x(x-2)$.\\
$f''(3)=36>0$: confirma el \textbf{mínim} en $x=3$.\\
$f''(0)=0$: el criteri \textbf{no decideix}. Cal tornar al signe de $f'$, que no canvia en
$x=0$, i per tant allà no hi ha extrem.
\end{solucio}
\end{nomesllarg}

\apartat[1,25]{0,75}
Demostra que la funció $h(x)=2x+\sin x$ és creixent a tot $\mathbb{R}$.

\begin{solucio}
$h'(x)=2+\cos x$. Com que $-1\le\cos x\le1$, es compleix $h'(x)\ge2-1=1>0$ per a tot $x$.\\
La derivada és sempre positiva, i per tant $h$ \textbf{creix a tot $\mathbb{R}$}.
\end{solucio}

\end{apartats}
